\section{Fixed-tube history graphs and mixed parameter jets}
\label{app:compact-graph-jets}
\label{app:compact-graph}

This appendix proves the fixed-tube graph result used in
\cref{sec:history-graph}.  The point requiring some care is that the stable
generator is \(A/\del\), while delayed evaluation is performed along the
flow of the unknown reduced vector field.  A fixed-parameter center-manifold
theorem therefore does not give estimates that remain regular at
\(\del=0\).  We instead use a derivative-free Lyapunov--Perron transform.
Its zeroth-order contraction is uniform at the singular limit, and its
mixed derivatives close on a finite triangular scale of Banach spaces.

\subsection{Prepared form of the fold chart}
\label{app:compact-graph-setting}

Recall the shifts
\begin{equation}
\label{eq:app-stable-shifts}
 \widehat Z=Z+\frac{\alpha}{2}X^2,
 \qquad
 \widehat W=W+\frac{\alpha}{2d_w}X^2,
 \qquad
 u=(X,Y)^T,
 \qquad
 h=(\widehat Z,\widehat W)^T.
\end{equation}
On a fixed compact flow tube, the exact equations obtained from
\cref{eq:blowup-scaling} have the form
\begin{equation}
\label{eq:app-abstract-chart}
\begin{aligned}
 u'(s)&=q_0(u(s))
 +\del F\!\left(u(s),h(s),
       (u(s-\theta_j),h(s-\theta_j))_{j=0,1};
       \del,\nu,\eta\right),\\
 \del h'(s)&=Ah(s)
 +\del G\!\left(u(s),h(s),
       (u(s-\theta_j),h(s-\theta_j))_{j=0,1};
       \del,\nu,\eta\right),
\end{aligned}
\end{equation}
where
\begin{equation}
\label{eq:app-q0-A}
 q_0(X,Y)=\vect{Y-\alpha X^2\\-X},
 \qquad
 A=\mat{-2&0\\1&-d_w}.
\end{equation}
The cancellations leading to the external factors \(\del\) in
\cref{eq:app-abstract-chart} are exact.  Since \(d_w>0\), \(A\) is Hurwitz;
we fix constants \(M_A\geq1\) and \(\beta>0\) such that
\begin{equation}
\label{eq:app-stable-semigroup}
 \norm{\e^{Ar}}\leq M_A\e^{-\beta r},
 \qquad r\geq0.
\end{equation}
This estimate also covers the non-diagonalizable case \(d_w=2\).

Choose a smooth preparation which agrees with the polynomial chart on a
neighborhood of the fixed tube.  Outside that neighborhood we require the
prepared \(q_0,F,G\) to belong to \(C_b^{12}\), including their state and
parameter derivatives, and require \(q_0\) to generate a complete
two-sided flow.  We also assume that the prepared data extend to an open
neighborhood of
\begin{equation}
\label{eq:app-parameter-box-star}
 [-\del_*,\del_*]\times I_\nu\times[-\eta_*,\eta_*],
\end{equation}
where \(I_\nu\Subset\R\) is fixed.  Such a preparation is obtained by a
standard compactly supported cutoff.  The extension to negative \(\del\)
will be used only for Taylor's theorem; the RFDE interpretation is made for
\(\del>0\).

The following proposition is the fixed-tube statement needed in the main
text.  Its constants may depend on the chosen preparation and on
\(I_\nu\), but not on \((\del,\nu,\eta)\) in the smaller parameter box.

\begin{proposition}[Fixed-tube special-flow graph]
\label{prop:app-fixed-tube-graph}
There are \(0<\del_0<\del_*\), \(0<\eta_0<\eta_*\), and \(C>0\) such that
the following assertions hold for
\[
 (\del,\nu,\eta)\in
 \Lambda_0:=[-\del_0,\del_0]\times I_\nu\times[-\eta_0,\eta_0].
\]
There is a unique pair \(Z=(Q,H)\) in a fixed \(C_b^1\) neighborhood of
\((q_0,0)\) satisfying the graph transform in
\cref{eq:app-graph-transform} below, and
\begin{equation}
\label{eq:app-C1-smallness}
 \norm{Q-q_0}_{C_b^1}+\norm{H}_{C_b^1}\leq C\abs{\del}.
\end{equation}
For \(\del>0\), this pair satisfies the invariance equations
\begin{equation}
\label{eq:app-invariance-equations}
\begin{aligned}
 Q(u)&=q_0(u)+\del F(\cE_{Q,H}(u);\del,\nu,\eta),\\
 \del DH(u)Q(u)&=AH(u)
       +\del G(\cE_{Q,H}(u);\del,\nu,\eta),
\end{aligned}
\end{equation}
where \(\cE_{Q,H}\) is defined in \cref{eq:app-history-evaluation}.
Moreover,
\begin{equation}
\label{eq:app-mixed-jet-bound}
 \max_{\substack{0\leq b\leq3\\0\leq i\leq1\\0\leq c\leq2}}
 \norm{\partial_\del^b\partial_\nu^i\partial_\eta^c(Q,H)}
       _{C_b^{\,9-b-i-c}}
 \leq C.
\end{equation}
In particular, \(Q,H\) are \(C^3\) in the state variables, three times
differentiable in \(\del\), once differentiable in \(\nu\), and twice
differentiable in \(\eta\), with all the displayed rectangular mixed
derivatives continuous also at \(\del=0\).

Let \(\Phi_Q^s\) be the complete flow of \(Q\) and put
\begin{equation}
\label{eq:app-history-embedding}
 \iota_{\del,\nu,\eta}(u)(\vartheta)
 =\bigl(\Phi_Q^\vartheta u,
        H(\Phi_Q^\vartheta u)\bigr),
 \qquad -\theta_1\leq\vartheta\leq0.
\end{equation}
For \(\del>0\), this is a \(C^3\) injective embedding into
\(C([ -\theta_1,0],\R^4)\).  If \(\mathcal S_s^\sharp\) denotes the
semiflow of the prepared equation, then
\begin{equation}
\label{eq:app-semiconjugacy}
 \mathcal S_s^\sharp\iota_{\del,\nu,\eta}(u)
 =\iota_{\del,\nu,\eta}(\Phi_Q^s u),
 \qquad s\geq0.
\end{equation}
Whenever the history segment on the right remains in the region where the
preparation agrees with the exact fold chart, it is also an exact history
segment of the unprepared equation.

Finally, in \(C_b^3\), uniformly together with one \(\nu\)-derivative and
two \(\eta\)-derivatives,
\begin{equation}
\label{eq:app-graph-expansion}
\begin{aligned}
 Q&=q_0+\del Q_1+\del^2Q_2+\del^3R_{3,Q},\\
 H&=\del H_1+\del^2H_2+\del^3R_{3,H},
\end{aligned}
\qquad
 \max_{\substack{0\leq i\leq1\\0\leq c\leq2}}
 \norm{\partial_\nu^i\partial_\eta^c(R_{3,Q},R_{3,H})}_{C_b^3}
 \leq C.
\end{equation}
If \(\phi_0^s\) is the flow of \(q_0\), then the first coefficients are
\begin{equation}
\label{eq:app-first-graph-coefficients}
 Q_1(u)=F(\cE_0(u);0,\nu,\eta),
 \qquad
 H_1(u)=-A^{-1}G(\cE_0(u);0,\nu,\eta),
\end{equation}
where
\begin{equation}
\label{eq:app-E0}
 \cE_0(u)=
 \left(u,0,
  (\phi_0^{-\theta_j}u,0)_{j=0,1}\right).
\end{equation}
\end{proposition}

The proposition constructs one invariant special-flow graph.  It does not
assert that every nearby RFDE solution lies on the graph, and no stable
foliation of the full history space is needed for the connection argument.

\subsection{A derivative-free Lyapunov--Perron transform}
\label{app:compact-LP-transform}

For a complete vector field \(Q\), let \(\Phi_Q^s\) denote its flow and set
\begin{equation}
\label{eq:app-history-evaluation}
 \cE_{Q,H}(u)=
 \left(u,H(u),
   \bigl(\Phi_Q^{-\theta_j}u,
         H(\Phi_Q^{-\theta_j}u)\bigr)_{j=0,1}\right).
\end{equation}
Define
\begin{equation}
\label{eq:app-graph-transform}
\begin{aligned}
 \cT_Q(Q,H)(u)
 &=q_0(u)+\del F(\cE_{Q,H}(u);\del,\nu,\eta),\\
 \cT_H(Q,H)(u)
 &=\del\int_0^\infty \e^{Ar}
 G\!\left(
   \cE_{Q,H}(\Phi_Q^{-\del r}u);
   \del,\nu,\eta\right)\dd r.
\end{aligned}
\end{equation}
The second formula contains no derivative of the unknown \(H\).  Its
integrand samples precisely the points
\begin{equation}
\label{eq:app-sampled-flow-points}
 \Phi_Q^{-\del r}u,
 \qquad
 \Phi_Q^{-(\theta_j+\del r)}u,
 \quad j=0,1.
\end{equation}
By \cref{eq:app-stable-semigroup}, the integral converges absolutely for
bounded \(G\).

Choose \(L>\Lip(q_0)\).  For fixed constants \(b_Q,b_H\), consider the
closed subset of bounded Lipschitz pairs given by
\begin{equation}
\label{eq:app-Lipschitz-ball}
\begin{aligned}
 \norm{Q-q_0}_\infty+\Lip(Q-q_0)&\leq b_Q\del_0,\\
 \norm{H}_\infty+\Lip(H)&\leq b_H\del_0,
 \qquad \Lip(Q)\leq L.
\end{aligned}
\end{equation}
It is complete in the uniform product norm.  Every \(Q\) in this set is
bounded and globally Lipschitz, hence has a complete two-sided flow.
For two such vector fields, Gronwall's inequality gives
\begin{equation}
\label{eq:app-flow-C0-difference}
 \sup_u\abs{\Phi_Q^t u-\Phi_{\widetilde Q}^t u}
 \leq \abs{t}\e^{L\abs{t}}
       \norm{Q-\widetilde Q}_\infty,
 \qquad t\in\R.
\end{equation}

The uniform bounds on the data immediately imply
\begin{equation}
\label{eq:app-transform-size}
 \norm{\cT_Q(Q,H)-q_0}_\infty\leq C\abs{\del},
 \qquad
 \norm{\cT_H(Q,H)}_\infty
 \leq \frac{M_A}{\beta}\norm{G}_\infty\abs{\del}.
\end{equation}
The Lipschitz estimate for a flow, together with
\cref{eq:app-sampled-flow-points}, yields
\begin{equation}
\label{eq:app-transform-Lipschitz-size}
\begin{aligned}
 \Lip(\cT_Q-q_0)&\leq C\abs{\del},\\
 \Lip(\cT_H)&\leq
 C\abs{\del}\int_0^\infty
       \e^{-(\beta-L\abs{\del})r}\dd r.
\end{aligned}
\end{equation}
We first choose \(b_Q,b_H\) larger than the constants in these estimates
and then decrease \(\del_0\) so that \(L\del_0<\beta/2\).  The transform
then maps \cref{eq:app-Lipschitz-ball} into itself.

For the difference of two pairs, each occurrence of a different flow in
\cref{eq:app-history-evaluation} is bounded by
\cref{eq:app-flow-C0-difference}.  In the integral component the resulting
majorant is a constant multiple of
\[
 \abs{\del}\e^{-\beta r}
 (1+\theta_1+\abs{\del}r)
 \e^{L(\theta_1+\abs{\del}r)}.
\]
It is integrable uniformly for \(\abs{\del}\leq\del_0\).  Hence
\begin{equation}
\label{eq:app-C0-contraction}
 \norm{\cT(Z)-\cT(\widetilde Z)}_{C^0}
 \leq C\del_0\norm{Z-\widetilde Z}_{C^0}.
\end{equation}
After one more decrease of \(\del_0\), the right-hand coefficient is a
number \(\kappa_0<1\).  Banach's fixed-point theorem gives a unique fixed
point in \cref{eq:app-Lipschitz-ball}.  Applying
\cref{eq:app-transform-size,eq:app-transform-Lipschitz-size} at that fixed
point gives the sharper slice estimate \cref{eq:app-C1-smallness}.

We next verify that the fixed point gives exact histories.  Let
\(u(s)=\Phi_Q^s u_*\), \(h(s)=H(u(s))\), and abbreviate the \(G\)-forcing
along this curve by \(g(s)\).  For \(\del>0\), the second fixed-point
equation is equivalent, after the change of variables
\(\tau=s-\del r\), to
\begin{equation}
\label{eq:app-stable-convolution}
 h(s)=\int_{-\infty}^s
       \e^{A(s-\tau)/\del}g(\tau)\dd\tau.
\end{equation}
The bounded continuous convolution is differentiable in \(s\), and
\[
 h'(s)=\del^{-1}Ah(s)+g(s).
\]
Together with the first fixed-point equation, this proves that
\((u(s),h(s))\) solves \cref{eq:app-abstract-chart}.  Conversely, a bounded
\(C^1\) pair in \cref{eq:app-Lipschitz-ball} satisfying
\cref{eq:app-invariance-equations} can be integrated along the complete
backward \(Q\)-orbit.  The boundary term is killed by
\cref{eq:app-stable-semigroup}; letting the backward time tend to infinity
recovers \cref{eq:app-graph-transform}.  Thus the fixed point is also unique
among bounded \(C^1\) solutions of the invariance equations in the stated
neighborhood.

The history identity \cref{eq:app-semiconjugacy} follows by taking segments
of this complete solution.  Evaluation at \(\vartheta=0\), followed by
projection to the \(u\)-coordinates, is a continuous left inverse of
\(\iota_{\del,\nu,\eta}\).  It follows that the history map is injective and
a homeomorphism onto its image.  Once the spatial regularity proved below
is available, it is \(C^3\); its differential is injective because the
present \(u\)-component of that differential is the identity.  This proves
all embedding and semiconjugacy assertions in
\cref{prop:app-fixed-tube-graph}.

\subsection{Triangular flow jets}
\label{app:compact-flow-jets}
\label{app:mixed-jets}

It remains to establish parameter regularity without applying a same-order
implicit-function theorem.  For a parameter-dependent map \(f\), write
\begin{equation}
\label{eq:app-jet-notation}
 J_{a,b,i,c}f=D_u^a\partial_\del^b\partial_\nu^i\partial_\eta^c f,
 \qquad
 |(a,b,i,c)|=a+b+i+c,
 \qquad
 \wp(a,b,i,c)=b+i+c.
\end{equation}
Set
\begin{equation}
\label{eq:app-triangular-index-set}
 \cI_k=
 \left\{(a,b,i,c)\in\mathbb N_0^4:
 b\leq3,\ i\leq1,\ c\leq2,\ a+b+i+c\leq k\right\}.
\end{equation}
The desired rectangular family
\(0\leq a\leq3\), \(0\leq b\leq3\), \(0\leq i\leq1\), and
\(0\leq c\leq2\) is contained in
\(\cI_9\).  We shall construct bounded jets through \(\cI_{10}\) and prove
convergence through \(\cI_9\); the last grade supplies the spatial reserve
needed when two flow jets are compared.

Order the indices first by total grade \(n=a+b+i+c\), and at equal total
grade by increasing parameter grade \(p=b+i+c\).  Indices with the same
\((n,p)\) form a block \(\cB_{n,p}\).  Thus, at a fixed total grade, jets
with more spatial derivatives are available before jets with more parameter
derivatives.

For \(\gamma=(a,b,i,c)\), let
\begin{equation}
\label{eq:app-jet-spaces}
 \cY_\gamma^e
 =C_b\!\left(\R^2,
     \mathcal L_{\rm sym}^a(\R^2;\R^e)\right),
 \qquad
 \widehat{\cY}_\gamma^e
 =C(\Lambda_0;\cY_\gamma^e),
\end{equation}
with the supremum norm over state and parameters.  The common fiber for a
block is
\begin{equation}
\label{eq:app-common-fiber}
 \mathfrak F_{n,p}
 =\prod_{\gamma\in\cB_{n,p}}
   \left(\widehat{\cY}_\gamma^2
         \times\widehat{\cY}_\gamma^2\right),
\end{equation}
where the two factors correspond to \(Q\) and \(H\).  This fiber is fixed
throughout the iteration.

For \(r\geq0\), introduce
\begin{equation}
\label{eq:app-flow-compositions}
 \Psi_{j,r}[Q](u)=\Phi_Q^{-(\theta_j+\del r)}u,
 \quad j=0,1,
 \qquad
 \Psi_{*,r}[Q](u)=\Phi_Q^{-\del r}u.
\end{equation}

\begin{lemma}[Triangular mixed jets of the flow]
\label{lem:app-triangular-flow-jets}
Fix \(k\leq10\) and suppose that a family \(Q\) has a common \(C_b^1\)
bound and bounded jets in \(\cI_k\).  For every
\(\gamma\in\cI_k\) there are constants \(C_\gamma,\Gamma_\gamma>0\)
and an integer \(M_\gamma\) such that
\begin{equation}
\label{eq:app-flow-jet-bound}
 \norm{J_\gamma\Psi_{j,r}[Q]}_\infty
 \leq C_\gamma(1+r^{M_\gamma})
 \e^{\Gamma_\gamma(\theta_1+\abs{\del}r)},
 \qquad j\in\{*,0,1\}.
\end{equation}
Suppose, in addition, that \(Q,\widetilde Q\) have common bounds through
the next spatial grade.  If \(d_{<\gamma}(Q,\widetilde Q)\) denotes the
maximum difference of all jet blocks preceding the block of \(\gamma\),
then
\begin{align}
\label{eq:app-flow-jet-difference}
 &\norm{J_\gamma\Psi_{j,r}[Q]
       -J_\gamma\Psi_{j,r}[\widetilde Q]}_\infty\notag\\
 &\quad\leq C_\gamma(1+r^{M_\gamma})
 \e^{\Gamma_\gamma(\theta_1+\abs{\del}r)}
 \left(
  \norm{J_\gamma Q-J_\gamma\widetilde Q}_\infty
  +d_{<\gamma}(Q,\widetilde Q)
 \right).
\end{align}
The same statements hold for \(H\circ\Psi_{j,r}[Q]\), with the current
block of \((Q,H)\) on the right-hand side.
\end{lemma}

\begin{proof}
First regard the terminal time \(t\) as independent of the parameters.
The flow solves
\[
 \partial_t\Phi_Q^t(u)=Q(\Phi_Q^t(u)).
\]
Its first spatial variation satisfies
\begin{equation}
\label{eq:app-first-variation}
 \partial_tD_u\Phi_Q^t
 =DQ(\Phi_Q^t)D_u\Phi_Q^t,
\end{equation}
and hence is bounded by \(\e^{L\abs t}\).  At spatial order \(a\geq2\),
the differentiated equation has the form
\begin{align}
\label{eq:app-higher-variation}
 \partial_tD_u^a\Phi_Q^t
 ={}&DQ(\Phi_Q^t)D_u^a\Phi_Q^t\notag\\
 &+D^aQ(\Phi_Q^t)
   [D_u\Phi_Q^t,\ldots,D_u\Phi_Q^t]
   +\mathcal R_a.
\end{align}
Every term in \(\mathcal R_a\) contains only lower spatial derivatives.
Induction and variation of constants give
\cref{eq:app-flow-jet-bound} for pure spatial jets.

For any parameter \(\lambda\in\{\del,\nu,\eta\}\),
\begin{equation}
\label{eq:app-parameter-variation}
 \partial_t\partial_\lambda\Phi_Q^t
 =DQ(\Phi_Q^t)\partial_\lambda\Phi_Q^t
  +(\partial_\lambda Q)(\Phi_Q^t).
\end{equation}
After arbitrary mixed differentiation, the equation for
\(J_\gamma\Phi_Q^t\) is linear in that jet.  In the inhomogeneity, the
occurrence of \(J_\gamma Q\) is linear.  Every other vector-field jet of
the same total grade has smaller parameter grade: this happens when a
parameter derivative strikes the argument \(\Phi_Q^t\), creating an
additional state slot on \(Q\).  Products containing two positive-grade
jets contain only factors of smaller total grade.  This is precisely the
block ordering introduced above.  Induction over those blocks and
variation of constants prove the fixed-time estimate.

We now substitute \(t=-(\theta_j+\del r)\).  Every \(\del\)-derivative
which strikes the terminal time produces a factor \(r\) and a time
derivative of the flow.  The latter is again expressed through the flow
equation and already controlled jets.  This changes only the polynomial
factor \(1+r^{M_\gamma}\), not the triangular highest-block structure.
Subtracting the two variational systems proves
\cref{eq:app-flow-jet-difference}; one additional spatial derivative
controls the evaluation of a highest derivative of \(Q\) at two different
orbits.  The ordinary Faà di Bruno formula applied to
\(H\circ\Psi_{j,r}[Q]\) gives the last assertion.
\end{proof}

The exponential factor in \cref{eq:app-flow-jet-bound} is harmless in the
stable integral.  Since only finitely many jets occur, we may decrease
\(\del_0\) so that
\begin{equation}
\label{eq:app-integrability-choice}
 \Gamma_{10}\del_0<\frac{\beta}{2},
\end{equation}
where \(\Gamma_{10}\) is the maximum of the exponents required through
grade ten, including the reserve used to prove convergence through grade
nine.  Every differentiated integrand is then dominated by
\begin{equation}
\label{eq:app-integrable-majorant}
 C(1+r^M)\e^{-\beta r/2},
\end{equation}
uniformly on \(\Lambda_0\).  This justifies differentiation under the
integral for every finite graph-transform iterate.

\subsection{Invariant jet fibers and convergence}
\label{app:compact-fiber-contraction}

Start with \(Z_0=(q_0,0)\) and set \(Z_{k+1}=\cT(Z_k)\).  Each finite
iterate is smooth.  Denote its block of jets by
\begin{equation}
\label{eq:app-block-jet}
 \mathbf J_{n,p}Z_k
 =(J_\gamma Z_k)_{\gamma\in\cB_{n,p}}
 \in\mathfrak F_{n,p},
\end{equation}
and let
\begin{equation}
\label{eq:app-lower-block-distance}
 d_{<n,p}(Z,\widetilde Z)
 =\max_{\cB_{n',p'}\prec\cB_{n,p}}
   \max_{\gamma\in\cB_{n',p'}}
   \norm{J_\gamma Z-J_\gamma\widetilde Z}_\infty,
\end{equation}
with an empty maximum equal to zero.

The next lemma isolates the highest block in the chain rule.  It is the
finite-dimensional-in-order substitute for a same-order differentiability
claim about the entire graph transform.

\begin{lemma}[Highest-block recurrence]
\label{lem:app-highest-block-recurrence}
Assume that all blocks preceding \(\cB_{n,p}\subset\cI_{10}\) range in fixed
bounded subsets of their common fibers.  Except for the first pure spatial
block \(\cB_{1,0}\), differentiation of \cref{eq:app-graph-transform} has
the affine form
\begin{equation}
\label{eq:app-affine-fiber-map}
 \mathbf J_{n,p}\cT(Z)
 =\mathbf b_{n,p}(\mathbf J_{<n,p}Z)
  +\del\,\mathscr L_{n,p}(\mathbf J_{<n,p}Z)
       [\mathbf J_{n,p}Z],
\end{equation}
where \(\mathscr L_{n,p}\) is a bounded linear map on
\(\mathfrak F_{n,p}\).  On the stated bounded sets,
\begin{equation}
\label{eq:app-fiber-operator-bound}
 \norm{\mathscr L_{n,p}}\leq C_{n,p}.
\end{equation}
For any two actual smooth jet families,
\begin{equation}
\label{eq:app-fiber-difference}
\begin{aligned}
 &\norm{\mathbf J_{n,p}\cT(Z)
       -\mathbf J_{n,p}\cT(\widetilde Z)}\\
 &\quad\leq C_{n,p}\abs{\del}
       \norm{\mathbf J_{n,p}Z-
              \mathbf J_{n,p}\widetilde Z}
       +C_{n,p}d_{<n,p}(Z,\widetilde Z).
\end{aligned}
\end{equation}
Estimate \cref{eq:app-fiber-difference} also holds for
\(\cB_{1,0}\).
\end{lemma}

\begin{proof}
Faà di Bruno's formula shows that, in a derivative of total grade \(n\),
a factor carrying the full grade can occur only once.  If a parameter
derivative enters a flow argument, the vector-field jet left outside has
the same total grade but smaller parameter grade and is therefore in a
preceding block.  The same triangular statement for the implicitly defined
flow jets is \cref{lem:app-triangular-flow-jets}.

For clarity, \(\mathscr L_{n,p}\) can be defined without invoking a
derivative of the graph transform on a same-order Banach space.  Replace
the current jet entries by an arbitrary
\(U\in\mathfrak F_{n,p}\).  In
\cref{eq:app-first-variation,eq:app-higher-variation,eq:app-parameter-variation},
retain the terms linear in \(U\) and solve the resulting finite family of
linear variational equations with zero initial values.  Substitute these
top flow jets, together with the \(H\)-entries of \(U\), into the terms
linear in the current block in the Faà di Bruno expansions of \(F\) and
\(G\).  Applying the semigroup integral to the latter defines
\(\mathscr L_{n,p}U\); setting \(U=0\) leaves
\(\mathbf b_{n,p}\).  The estimates in
\cref{lem:app-triangular-flow-jets,eq:app-integrable-majorant} prove
\cref{eq:app-fiber-operator-bound}.

Both components of \cref{eq:app-graph-transform} carry an external
factor \(\del\).  If a \(\del\)-derivative removes that factor, then
\begin{equation}
\label{eq:app-Leibniz-external-delta}
 \partial_\del^b(\del\mathcal F)
 =\del\partial_\del^b\mathcal F
  +b\partial_\del^{b-1}\mathcal F,
\end{equation}
and the second term has lower total grade.  Thus every occurrence of the
current block retains the factor \(\del\).  Subtracting two affine formulas
and using \cref{eq:app-flow-jet-difference} gives
\cref{eq:app-fiber-difference}.

For \(\cB_{1,0}\), the first variational equation contains \(DQ\) in its
coefficient, so dependence on the current block need not be affine.
Nevertheless, the difference of two first variational equations, followed
by Gronwall's inequality, is uniformly Lipschitz on the invariant \(C_b^1\)
ball.  One reserved spatial derivative controls evaluation at different
orbits.  The external factor \(\del\) is still present, which proves
\cref{eq:app-fiber-difference} for this block as well.
\end{proof}

\begin{lemma}[Uniform invariant jet balls]
\label{lem:app-invariant-jet-balls}
After decreasing \(\del_0\), there are radii \(R_{n,p}\), independent of
the iteration number and of the point in \(\Lambda_0\), such that
\begin{equation}
\label{eq:app-invariant-jet-balls}
 \norm{\mathbf J_{n,p}Z_k}_{\mathfrak F_{n,p}}
 \leq R_{n,p}
\end{equation}
for every block in \(\cI_{10}\) and every \(k\geq0\).
\end{lemma}

\begin{proof}
The order-zero block is bounded by the Lipschitz contraction ball, and the
first pure spatial block is bounded by the invariant \(C_b^1\) estimate.
Suppose the radii have been fixed for every preceding block.  The preceding
lemma, \cref{eq:app-Leibniz-external-delta}, and the fixed derivatives of
\(q_0\) give
\begin{equation}
\label{eq:app-invariant-recurrence}
 \norm{\mathbf J_{n,p}Z_{k+1}}
 \leq C_{n,p}\abs{\del}
       \norm{\mathbf J_{n,p}Z_k}+B_{n,p},
\end{equation}
where \(B_{n,p}\) depends only on the already chosen lower-block radii and
the prepared data.  All terms in the stable component are finite by
\cref{eq:app-integrable-majorant}.  Choose \(\del_0\) so that
\(C_{n,p}\del_0\leq1/2\) for the finitely many blocks, and then take
\[
 R_{n,p}\geq
 2B_{n,p}+\norm{\mathbf J_{n,p}Z_0}.
\]
Induction first over the ordered blocks and then over \(k\) proves
\cref{eq:app-invariant-jet-balls}.
\end{proof}

\begin{lemma}[Convergence of the mixed jets]
\label{lem:app-mixed-jet-convergence}
For every block in \(\cI_9\), the sequence
\(\mathbf J_{n,p}Z_k\) converges uniformly on
\(\R^2\times\Lambda_0\).
\end{lemma}

\begin{proof}
The order-zero block converges geometrically by
\cref{eq:app-C0-contraction}.  Suppose that all preceding blocks converge
geometrically.  The bound through \(\cI_{10}\) from
\cref{lem:app-invariant-jet-balls} supplies the reserved derivative required
in \cref{eq:app-flow-jet-difference}.  Applying
\cref{eq:app-fiber-difference} to two consecutive iterates gives
\begin{align}
\label{eq:app-fiber-increment-recurrence}
 \norm{\mathbf J_{n,p}Z_{k+1}-\mathbf J_{n,p}Z_k}
 \leq{}&\kappa_{n,p}
 \norm{\mathbf J_{n,p}Z_k-\mathbf J_{n,p}Z_{k-1}}\notag\\
 &+C_{n,p}d_{<n,p}(Z_k,Z_{k-1}),
\end{align}
where \(\kappa_{n,p}\leq C_{n,p}\del_0<1\).  By induction the last term is
bounded by \(Cq^k\) for some \(q<1\).  The scalar recurrence
\[
 a_{k+1}\leq\kappa a_k+Cq^k
\]
has a geometric bound at every rate strictly larger than
\(\max\{\kappa,q\}\).  Thus the increments in the current fiber are
summable.  The finite block induction proves the claim.
\end{proof}

Let \(Z=(Q,H)\) be the \(C^0\) fixed point constructed above.  The preceding
lemma supplies uniformly convergent candidates for all its jets in
\(\cI_9\).  We record why they are actual derivatives.  If
\(f_k\to f\) and \(\partial_\lambda f_k\to g\) uniformly in one common
fiber, then on every coordinate segment contained in the open parameter
neighborhood,
\begin{equation}
\label{eq:app-strong-jet-closure}
 f_k(u,\lambda+t)-f_k(u,\lambda)
 =\int_0^t\partial_\lambda f_k(u,\lambda+r)\dd r.
\end{equation}
Passing to the limit gives the same identity for \(f,g\), hence
\(\partial_\lambda f=g\).  The identical line-segment argument in \(u\)
identifies spatial derivatives.  Iteration in the block order identifies
every limiting jet as \(J_{a,b,i,c}Z\).  At the boundary of \(\Lambda_0\),
the derivatives are the restrictions of those constructed on the open
parameter neighborhood.  Consequently
\[
 \max_{\substack{0\leq b\leq3,\ 0\leq i\leq1\\0\leq c\leq2}}
 \norm{\partial_\del^b\partial_\nu^i\partial_\eta^cZ}
       _{C_b^{\,9-b-i-c}}<\infty,
\]
which is \cref{eq:app-mixed-jet-bound}.  This is strong convergence of
derivatives in fixed Banach fibers; no compactness or weak-limit argument is
used.

\subsection{Taylor coefficients and the uniform remainder}
\label{app:compact-Taylor-remainder}

At \(\del=0\), the transform \cref{eq:app-graph-transform} is the constant
map \(Z=(q_0,0)\), independently of \((\nu,\eta)\).  Define
\begin{equation}
\label{eq:app-Taylor-coefficients}
 (Q_j,H_j)=\frac1{j!}
 \partial_\del^j(Q,H)\big|_{\del=0},
 \qquad j=1,2.
\end{equation}
By \cref{eq:app-mixed-jet-bound}, Taylor's theorem applies with values in
\(C_b^3(\R^2,\R^4)\), also after at most one \(\nu\)-derivative and two
\(\eta\)-derivatives.  For
\(Z=(Q,H)\), \(Z_0=(q_0,0)\), and \(Z_j=(Q_j,H_j)\), the remainder is
\begin{equation}
\label{eq:app-integral-Taylor-remainder}
 (R_{3,Q},R_{3,H})
 =\frac12\int_0^1(1-t)^2
   \partial_\del^3Z(t\del,\nu,\eta)\dd t.
\end{equation}
Applying \(\partial_\nu^i\partial_\eta^c\),
\(0\leq i\leq1\), \(0\leq c\leq2\), and using
\cref{eq:app-mixed-jet-bound} proves the uniform remainder in
\cref{eq:app-graph-expansion}.

Finally, differentiate \cref{eq:app-graph-transform} once at \(\del=0\).
All derivatives of the inner arguments are multiplied by the vanishing
external factor.  Therefore
\[
 Q_1=F(\cE_0;0,\nu,\eta),
 \qquad
 H_1=\left(\int_0^\infty\e^{Ar}\dd r\right)
       G(\cE_0;0,\nu,\eta).
\]
Since \(A\) is Hurwitz,
\(\int_0^\infty\e^{Ar}\dd r=-A^{-1}\), which gives
\cref{eq:app-first-graph-coefficients}.  This completes the proof of
\cref{prop:app-fixed-tube-graph}. \hfill\(\square\)

\begin{remark}[Why one spatial reserve is necessary]
\label{rem:app-spatial-reserve}
The map \(Q\mapsto\Phi_Q^t\) loses one state derivative in a same-order
vector-field norm.  Indeed, if
\(V=D_Q\Phi_Q^t[\widehat Q]\), then
\[
 \dot V=DQ(\Phi_Q^t)V+\widehat Q(\Phi_Q^t),
\]
and three spatial derivatives of this equation contain a term involving
\(D^4Q\).  Thus a \(C_b^3\)-to-\(C_b^3\) implicit-function argument would
request a derivative absent from its domain.  The triangular scale above
keeps that derivative in the bounded grade \(\cI_{10}\) while proving
convergence only through \(\cI_9\); no infinite-dimensional smoothing or
Nash--Moser argument is required for the finite jet family used here.
\end{remark}
